%%%%%%%%%%%%%%%%%%%%%%%%%%%%%%%%%%%%%%%%%%%%%%%%%%%%%%%%%%%%%%%%%%%%%%
% Lớp 11 | Dãy số. Cấp số cộng. Cấp số nhân | Cấp số cộng
% Dạng 1D2?2-6: Tính tổng của cấp số cộng
% Số câu: 63  (sinh tự động bởi tach_cau.py; từ nay thầy sửa trực tiếp file này)
%%%%%%%%%%%%%%%%%%%%%%%%%%%%%%%%%%%%%%%%%%%%%%%%%%%%%%%%%%%%%%%%%%%%%%

%%%%%%%%%%%%%%%% TRẮC NGHIỆM 4 PHƯƠNG ÁN %%%%%%%%%%%%%%%%

% nguồn: Ok_Toan11_TracNghiem_Lan2
\begin{ex}%[1D2N2-6]%[ma:ab2a796f]
	Cho cấp số cộng có số hạng đầu $u_1=-\dfrac{1}{2}$  và công sai $d=\dfrac{1}{2}$. Năm số hạng liên tiếp đầu tiên lần lượt là
	\choice 
		{$-\dfrac{1}{2}; \, 0; \, 1; \, \dfrac{1}{2};  \, 1$}
		{$-\dfrac{1}{2}; \, 0; \, \dfrac{1}{2}; \, 0;  \, \dfrac{1}{2}$}
		{$\dfrac{1}{2}; \, 1; \, \dfrac{3}{2}; \ 2; \, \dfrac{5}{2}$} 
		{\True $-\dfrac{1}{2}; \, 0; \,  \dfrac{1}{2}; \, 1; \, \dfrac{3}{2}$}
	 \loigiai{
		Sử dụng công thức $u_n=u_1+(n-1) d$ với $u_1=-\dfrac{1}{2}, \  d=\dfrac{1}{2}$. Ta có
		\[ 
		\begin{aligned}
			 u_1&=-\dfrac{1}{2}; \\
			 u_2&=u_1+d=0; \\
			u_3&=u_2+d=\dfrac{1}{2};  \\
			u_4&=u_3+d=1;  \\
			u_5&=u_4+d=\dfrac{3}{2}. 
		\end{aligned}
		 \]
		Vậy $5$ số hạng đầu tiên là $-\dfrac{1}{2}; \, 0; \,  \dfrac{1}{2}; \, 1; \, \dfrac{3}{2}$.
	 }
\end{ex}

% nguồn: Ok_Toan11_TracNghiem_Lan2
\begin{ex}%[1D2H2-6]%[ma:77c902c9]
	Cho  cấp số cộng $(u_n)$ thỏa $\heva{& u_2-u_3+u_5=10\\& u_4+u_6=26}$. Tính $S=u_1+u_4+u_7+\ldots+u_{2\,011}$
	 \choice
		{\True  $S=2\,023\,736$}
		 {$S=2\,023\,563$}
		 {$S=6\,730\,444$}
		{$S=6\,734\,134$} 
	\loigiai{ 
		Ta có \begin{eqnarray*}
		&& \heva{& u_2-u_3+u_5=10\\& u_4+u_6=26}\\ &\Leftrightarrow& \heva{&u_1+d-u_1-2d+u_1+4d=10\\  &u_1+3d+u_1+5d=26} \\
		&  \Leftrightarrow& \heva{&u_1+3d=10\\ &2u_1+8d=26}\\
		&\Leftrightarrow& \heva{&u_1=1\\  &d=3.}
	\end{eqnarray*}
	Suy ra  $u_4=10$; $u_7=19$; $u_{10}=28$.\\
	Ta có $u_1; u_4;  u_7; u_{10}; \ldots; u_{2\,011}$ là cấp số cộng có $\heva{& u_1=1\\ &d=9\\& n=671.}$ \\
	Do đó \[S=\dfrac{671}{2}(2\cdot 1+670\cdot  9)=2\,023\,736.\]
} 
\end{ex}

% nguồn: Ok_Toan11_TracNghiem_Lan2
\begin{ex}%[1D2H2-6]%[ma:cf14ffa0]
	Cho cấp số cộng $(u_n)$ có $u_{2\,013}+u_6=1\,000$.  Tổng $2\,018$ số hạng đầu tiên của cấp số cộng đó là
	\choice 
		 {\True $1\,009\,000$}
		 {$100\,800$}
		{$1\,008\,000$} 
		 {$100\,900$}
	\loigiai{Gọi $d$ là công sai  của cấp số cộng. Khi đó\\
		 $u_{2\,013}+u_6=1\,000\Leftrightarrow u_1+2\,012d+u_1+5d=1\,000 \Leftrightarrow 2u_1+2\,017d=1\,000$.\\
		Ta có $S_{2\,018}=2\,018u_1+\dfrac {2\,017\cdot 2\,018}{2}d=1\,009\cdot  \left(2u_1+2\,017d\right) =1\,009\,000$.
	 }
\end{ex}

% nguồn: Ok_Toan11_TracNghiem_Lan2
\begin{ex}%[1D2H2-6]%[ma:307adcd1]
	Cho cấp số cộng $(u_n)$ có  $u_5=-15$; $u_{20}=60$. Tổng $20$ số hạng đầu tiên của cấp số cộng là
	\choice 
		{\True $S_{20}=250$} 
		 {$S_{20}=200$}
		{$S_{20}=-200$} 
		 {$S_{20}=-25$}
	\loigiai{Ta  có
		 \allowdisplaybreaks
		\begin{eqnarray*}
			&&\heva{&u_5=-15\\&u_{20}=60}\\ 
			&\Leftrightarrow& \heva{& u_1+4d=-15\\&u_1+19d=60}\\ 
			&\Leftrightarrow&  \heva{&u_1=-35\\&d=5.}
		\end{eqnarray*}
		Vậy $S_{20}=\dfrac  {(u_1+u_{20})\cdot 20}{2}=250$.
	} 
\end{ex}

% nguồn: Ok_Toan11_TracNghiem_Lan2
\begin{ex}%[1D2H2-6]%[ma:96f0d408]
	Cho cấp số cộng $(u_n)$  với $u_n=3-2n$ thì $S_{60}$ bằng
	\choice 
		{$-6\,960$} 
		 {$-117$}
		{\True  $-3\,840$}
		 {$-116$}
	\loigiai{Ta có $u_{n+1}=1-2n$. Ta có $u_{n+1}-u_n=-2$, $\forall n\in \mathbb{N^*}$, suy ra  $(u_n)$ là cấp số cộng có $ u_1=1$ và công sai $ d=-2$.\\
		Vậy $S_{60}=\dfrac {60}{2}\left(2u_1+59d\right)=-3\,840$. 
	 }
\end{ex}

% nguồn: Ok_Toan11_TracNghiem_Lan2
\begin{ex}%[1D2H2-6]%[ma:f61dbb11]
	 Trong sân vận động có tất cả $30$ dãy ghế, dãy đầu tiên có $15$ ghế. Các dãy sau, mỗi dãy nhiều hơn dãy ngay trước nó $4$ ghế. Hỏi sân vận động có tất cả bao nhiêu ghế?
	 \choice
		{$1\,740$} 
		 {$2\,250$}
		{$4\,380$} 
		 {\True $2\,190$}
	\loigiai{Số ghế trong mỗi dãy của sân vận động lập  thành một cấp số cộng có $u_1=15$ và $d=4$.\\
		Vậy tổng tất cả các ghế của sân vận động là tổng $30$ số hạng đầu của cấp số cộng trên,  do đó áp dụng công thức $S_n=n\cdot u_1+\dfrac {n(n-1)d}{2}$.\\
		Ta có $S_{30}=30\cdot  15+\dfrac {30(30-1)4}{2}=2\,190$.\\
		Vậy sân vận  động có tất cả $2\,190$ ghế.
	 }
\end{ex}

% nguồn: Ok_Toan11_TracNghiem_Lan2
\begin{ex}%[1D2H2-6]%[ma:dec7b287]
	Biết $\lim\limits_{x\to 3}\dfrac{x^2+bx+c}{x-3}=8$  ($b$, $c\in\mathbb{R}$). Tính $P=b+c$.
	\choice 
		{\True$P=-13$} 
		{$P=-11$} 
		{$P=5$} 
		{$P=-12$} 
	 \loigiai{
		Vì $\lim\limits_{x\to 3}\dfrac{x^2+bx+c}{x-3}=8$ là hữu  hạn nên tam thức $x^2+bx+c$ có nghiệm $x=3$ khi và chỉ khi $3b+c+9=0\Leftrightarrow c=-9-3b$.
		Khi đó 
		\begin{eqnarray*}
			& &\lim\limits_{x\to 3}\dfrac{x^2+bx+c}{x-3}\\ 
			&=&\lim\limits_{x\to  3}\dfrac{x^2+bx-9-3b}{x-3}\\
			&=&\lim\limits_{x\to 3} \dfrac{(x-3)(x+3+b)}{x-3}\\ 
			&=&\lim\limits_{x\to 3}  (x+3+b)=6+b.
		\end{eqnarray*}
		Suy ra  $6+b=8\Leftrightarrow b=2$ suy ra $c=-15$.\\
		Vậy  $P=b+c=-13$.
	 }
\end{ex}

% nguồn: Ok_Toan11_TracNghiem_Lan2
\begin{ex}%[1D2H2-6]%[ma:f2d299e4]
	Cho cấp số cộng $(u_n)$ biết $u_3=6$, $u_8=16$. Tính công sai $d$ và tổng của $10$ số hạng đầu tiên. 
	 \choice
		{$d=2$;  $S_{10}=100$}
		{$d=1$; $S_{10}=80$} 
		 {$d=2$; $S_{10}=120$}
		{\True  $d=2$; $S_{10}=110$}
	 \loigiai{Ta có
		\allowdisplaybreaks 
		\begin{eqnarray*}
			 &&\heva{&u_3=6\\&u_8=16}\\
			&\Leftrightarrow& \heva{&u_1+2d=6\\&u_1+7d=16}\\ 
			&\Leftrightarrow& \heva{&u_1=2\\&d=2.} 
		\end{eqnarray*}
		Suy ra $S_{10}=10\cdot u_1+\dfrac {10\left(10-1\right)}{2}\cdot  d=10\cdot 2+\dfrac {10\left(10-1\right)}{2}\cdot 2=110$.
	 }
\end{ex}

% nguồn: Ok_Toan11_TracNghiem_Lan2
\begin{ex}%[1D2H2-6]%[ma:4b3f695b]
	Cho cấp số cộng $(u_n)$ thỏa mãn $\heva{&u_1+u_4=8\\& u_3-u_2=2}$. Tính tổng $10$ số hạng đầu của cấp số  cộng trên.
	 \choice
		{\True  $100$}
		{$110$} 
		 {$10$}
		 {$90$}
	\loigiai{ 
		Gọi cấp số cộng có  công sai là $d$ ta có $u_2=u_1+d$; $u_3=u_1+2d$; $u_4=u_1+3d$.\\
		Khi  đó
		\begin{eqnarray*}
			&& \heva{& u_1+u_4=8\\&  u_3-u_2=2 }\\
			&\Leftrightarrow& \heva{&2u_1+3d=8\\ &d=2}\\ 
			&\Leftrightarrow& \heva{&  u_1=1\\& d=2.}
		\end{eqnarray*}
		 Áp dụng công thức $S=nu_1+\dfrac{n(n-1)}{2}\cdot d$.\\
		Vậy tổng của $10$ số hạng đầu của cấp  số cộng là $S_{10}=10\cdot 1+\dfrac{10\cdot 9}{2}\cdot 2=100$.
	} 
\end{ex}

% nguồn: Ok_Toan11_TracNghiem_Lan2
\begin{ex}%[1D2H2-6]%[ma:9f4e66af]
	Một cấp số cộng có tổng của  $n$ số hạng đầu $S_n$ tính theo công thức $S_n=5n^2+3n, (n \in \mathbb{N}^*)$. Tìm số hạng đầu $u_1$ và công sai $d$ của cấp số cộng đó.
	 \choice
		{$u_1=-8$;  $d=10$}
		{$u_1=-8$; $d=-10$} 
		{\True $u_1=8$;  $d=10$}
		 {$u_1=8$; $d=-10$}
	\loigiai{ 
		Ta  có $u_1=S_1=8$; $u_2=S_2-S_1=18 \Rightarrow d=u_2-u_1=18-8=10$.
	} 
\end{ex}

% nguồn: Ok_Toan11_TracNghiem_Lan2
\begin{ex}%[1D2H2-6]%[ma:47f3032d]
	Cho cấp  số cộng $(u_n)$ và gọi $S_n$ là tổng $n$ số hạng đầu tiên của nó. Biết $S_7=77$ và $S_{12}=192$. Tìm số hạng tổng quát $u_n$ của cấp số cộng đó
	\choice 
		 {$u_n=5+4n$}
		{\True  $u_n=3+2n$}
		 {$u_n=2+3n$}
		{$u_n=4+5n$} 
	\loigiai{ 
		Giả sử cấp số cộng  có số hạng đầu là $u_1$ và công sai $d$.\\
		Ta có 
		\begin{eqnarray*}
			 &&\heva{&S_7=77\\ &S_{12}=192}\\
			&\Leftrightarrow& \heva{&7u_1+\dfrac{7\cdot 6\cdot d}{2}=77\\&12u_1+\dfrac{12\cdot 11\cdot d}{2}=192}\\ 
			 &\Leftrightarrow&\heva{&7u_1+21d=77\\ &12u_1+66d=192}\\
			&\Leftrightarrow&\heva{& u_1=5\\& d=2.} 
		\end{eqnarray*}
		Khi  đó $u_n=u_1+(n-1)d=5+2(n-1)=3+2n$.
	 }
\end{ex}

% nguồn: Ok_Toan11_TracNghiem_Lan2
\begin{ex}%[1D2H2-6]%[ma:44bb8997]
	Cho cấp số cộng $(u_n)$ biết  $u_3=6$; $u_8=16$. Tính công sai $d$ và tổng của $10$ số hạng đầu tiên.
	\choice 
		{$d=2; S_{10}=100$} 
		{$d=1;  S_{10}=80$}
		{$d=2; S_{10}=120$} 
		{\True $d=2; S_{10}=110$} 
	 \loigiai{
		Ta có 
		\begin{eqnarray*}
			&&\heva {&u_3=6\\  &u_8=16}\\
			& \Leftrightarrow& \heva {&u_1+2d=6\\  &u_1+7d=16}\\
			&\Leftrightarrow&  \heva {&u_1=2\\ &d=2} .\\
		\end{eqnarray*}
		\begin{eqnarray*}
			 S_{10}&=&10\cdot u_1+\dfrac{10(10-1)}{2}\cdot d\\
			 &=&10\cdot 2+\dfrac{10(10-1)}{2}\cdot 2\\
			&=&110. 
		\end{eqnarray*}
	} 
\end{ex}

% nguồn: Ok_Toan11_TracNghiem_Lan2
\begin{ex}%[1D2H2-6]%[ma:f2589fd4]
	Cho cấp số cộng $(u_n)$ với $u_n =  3 - 2n$ thì $S_{60}$ bằng
	 \choice
		{$-6\,960$} 
		{$  -117$}
		 {\True $-3\,480.$}
		{$  -116$}
	\loigiai{ 
		Ta có $u_1 = 3 - 2\cdot 1 = 1$.\\ 
		 Công sai $d = u_{n+1} - u_n = 3 - 2(n+1) - (3 - 2n) = 3 - 2n - 2 - 3 + 2n = -2$.\\
		 Tổng của $60$ số hạng đầu tiên là
		\begin{eqnarray*}
			S_{60} &=& \dfrac{60}{2}\left[2u_1 + (60-1)d\right]\\ &=& 30\left[2 \cdot 1 +  59 \cdot (-2)\right]\\
			&  = &30(2 - 118) = 30(-116)\\
			& =& -3\,480. 
		\end{eqnarray*}
	 }
\end{ex}

% nguồn: Ok_Toan11_TracNghiem_Lan2
\begin{ex}%[1D2H2-6]%[ma:c26bc168]
	Gọi $S_n$ là tổng $n$ số hạng đầu tiên trong cấp số cộng $(a_n)$. Biết $S_6=S_9$,  tỉ số $\dfrac{a_3}{a_5}$ bằng
	 \choice
		 {$\dfrac{9}{5}$}
		{$\dfrac{5}{9}$} 
		 {\True $\dfrac{5}{3}$}
		{$\dfrac{3}{5}$} 
	 \loigiai{
		 Ta có
		\begin{eqnarray*}
			&&S_6=S_9\\ 
			& \Leftrightarrow&  \dfrac{6(2a_1+5d)}{2}=\dfrac{9(2a_1+8d)}{2}\\ &\Leftrightarrow& a_1=-7d.
		\end{eqnarray*}
		Suy ra 
		\begin{eqnarray*}
			 \dfrac{a_3}{a_5}&=&\dfrac{a_1+2d}{a_1+4d}\\
			&=&\dfrac{-7d+2d}{-7d+4d}\\ 
			 &=&\dfrac{-5d}{-3d}\\
			 &=&\dfrac{5}{3}.
		\end{eqnarray*}
	} 
\end{ex}

% nguồn: Ok_Toan11_TracNghiem_Lan2
\begin{ex}%[1D2H2-6]%[ma:19b0521d]
	Cho cấp số cộng $(u_n)$ có $u_4=-12$; $u_{14}=18$. Tổng của  $16$ số hạng đầu tiên của cấp số cộng là
	\choice 
		{\True  $S=24$}
		{$S=-25$} 
		 {$S=-24$}
		{$S=26$} 
	\loigiai{ 
		 Ta có
		\begin{eqnarray*}
			&&\heva{&u_4=-12\\ &u_{14}=18}\\ 
			&\Leftrightarrow&\heva{&u_1+3d=-12\\& u_1+13d=18}\\ 
			&\Leftrightarrow&\heva{&u_1=-21\\  &d=3.}
		\end{eqnarray*}
		Tổng của $16$ số  hạng đầu tiên của cấp số cộng là $S_{16}=16\cdot(-21)+\dfrac{16\cdot 15}{2}\cdot 3=24$.
	} 
\end{ex}

% nguồn: Ok_Toan11_TracNghiem_Lan2
\begin{ex}%[1D2H2-6]%[ma:efeb8563]
	Cho cấp số cộng $(u_n)$  thỏa $\heva{&u_2-u_3+u_5=10\\&u_4+u_6=26}$. Tính $S=u_1+u_4+u_7+\cdots +u_{2\,011}$
	 \choice
		 {\True $S=2\,023\,736$}
		 {$S=2\,023\,563$}
		{$S=6\,730\,444$} 
		 {$S=6\,734\,134$}
	 \loigiai{Ta có
		 \allowdisplaybreaks
		\begin{eqnarray*}
			&&\heva{&u_2-u_3+u_5=10\\&u_4+u_6=26}\\ 
			 &\Leftrightarrow& \heva{& u_1+d-u_1-2d+u_1+4d=10\\&u_1+3d+u_1+5d=26}\\
			&\Leftrightarrow& \heva{&u_1+3d=10  \\&2u_1+8d=26}\\
			 &\Leftrightarrow& \heva{&u_1=1\\&d=3.}
		\end{eqnarray*}
		 Tính được $u_4=10$, $ u_7=19$, $ u_{10}=28$.\\
		Ta có $u_1$, $u_4$, $u_7$, $u_{10}$, ..., $ u_{2\,011}$ là cấp số cộng có  $\heva{&u_1=1\\&d=9\\&n=671.}$\\
		$S=\dfrac {671}{2}(2\cdot 1+670\cdot  9)=2\,023\,736$.
	 }
\end{ex}

% nguồn: Ok_Toan11_TracNghiem_Lan2
\begin{ex}%[1D2H2-6]%[ma:acf528ef]
	Cho cấp số cộng $(u_n)$ với số hạng đầu $u_1 = -6$ và công sai $d=4$. Tính tổng $S$  của $14$ số hạng đầu tiên của cấp số cộng đó.
	\choice 
		 {$46$}
		 {$308$}
		{$644$} 
		{\True $280$} 
	\loigiai{ 
		Tổng $14$ số hạng đầu tiên của  một cấp số cộng là $S_{14} = \dfrac{\left[2\cdot (-6)+(14-1)\cdot 4\right]\cdot 14}{2} = 280$.
	} 
\end{ex}

% nguồn: Ok_TracNghiemLop11_Lan1
\begin{ex}%[1D2H2-6]%[ma:2722904b]
	Cho  cấp số cộng $\left(u_n\right)$ có $u_2+u_8+u_9+u_{15}=100$. Khi đó $S_{16}$ bằng
	\choice 
		{$200$} 
		{$300$} 
		 {\True $400$}
		 {$100$}
	\loigiai{ 
		Gọi $d$ là công  sai của cấp số cộng.\\
		Ta có 
		\begin{eqnarray*}
			u_2+u_8+u_9+u_{15}=100&\Leftrightarrow& u_1+d+u_1+7d+u_1+8d+u_1+14d=100\\ 
			&\Leftrightarrow& 2u_1+15d=50. 
		\end{eqnarray*}
		Tổng $16$ số hạng đầu $S_{16}=16u_1+\dfrac{16\cdot 15}{2}d=8\left(2u_1+15d\right)=8\cdot  50=400$.
	 }
\end{ex}

% nguồn: Ok_TracNghiemLop11_Lan1
\begin{ex}%[1D2H2-6]%[ma:e9831b95]
	Cho cấp số cộng $\left(u_n \right)$ với $u_n=3-2n$  thì $S_{60}$ bằng
	 \choice
		 {$-6960$}
		{$-117$} 
		 {\True $-3840$}
		 {$-116$}
	\loigiai{ 
		 Ta có $u_{n+1}=1-2n$, Ta có $u_{n+1}-u_n=-2, \forall n\in \mathbb{N}^{*}$, suy ra $\left(u_n \right)$ là cấp số cộng có $u_1=1$ và công sai $d=-2$. \\
		Vậy  $S_{60}=\dfrac{60}{2} \left(2u_1+59d\right)=-3840$.
	 }
\end{ex}

% nguồn: Ok_TracNghiemLop11_Lan1
\begin{ex}%[1D2H2-6]%[ma:f26c7ec6]
	Cho cấp số cộng $\left(u_n \right)$ có $u_5=-15$, $u_{20}=60$. Tổng $S_{20}$  của $20$ số hạng đầu tiên của cấp số cộng là
	\choice 
		{$S_{20}=600$} 
		{$S_{20}=60$} 
		{\True $S_{20}=250$} 
		 {$S_{20}=500$}
	 \loigiai{
		Ta có $\heva{&{u_5=-15}  \\
		 &{u_{20}=60}
		&}$$\Leftrightarrow  \heva{&{u_1+4d=-15} \\
		&{u_1+19d=60} 
		&}$$\Leftrightarrow \heva{&{u_1=-35}  \\
		 &{d=5.}
		}$. 
		\[\Rightarrow  S_{20}=20u_1+\dfrac{20\cdot19}{2}\cdot d=20\cdot\left(-35\right)+\dfrac{20\cdot19}{2}\cdot5=250.
		\] 
	 }
\end{ex}

% nguồn: Ok_TracNghiemLop11_Lan1
\begin{ex}%[1D2H2-6]%[ma:6e385867]
	Một gia đình cần khoan một cái giếng để lấy nước. Họ thuê một đội khoan giếng nước đến để khoan giếng nước. Biết giá của mét khoan đầu tiên là $80\,000$ đồng, kể từ mét khoan thứ $2$ giá của mỗi  mét khoan tăng thêm $5\,000$ đồng so với giá của mét khoan trước đó. Biết cần phải khoan sâu xuống $50$ m mới có nước. Vậy hỏi phải trả bao nhiêu tiền để khoan cái giếng đó?
	 \choice
		{\True $10\,125\,000$  đồng}
		{$4\,000\,000$ đồng} 
		{$4\,245\,000$  đồng}
		{$5\,2500\,000$ đồng} 
	\loigiai{ 
		Giá tiền khoan mỗi mét (bắt đầu từ mét đầu tiên) lập thành  cấp số cộng $(u_n)$ có $u_1=80\,000$ và $d=5\,000$.\\
		Do cần khoan $50$ mét nên tổng số  tiền cần trả là
	 $$S_{50}=50\cdot 80\,0000 + \dfrac{50(50-1)}{2}\cdot 5\,000=10\,125\,000\,\, \text{(đồng)}.$$ }
\end{ex}

% nguồn: Ok_TracNghiemLop11_Lan1
\begin{ex}%[1D2H2-6]%[ma:11d8abb5]
	Cho cấp số cộng $\left(u_n\right)$ có số hạng đầu $u_1=-1$ và  công sai $d=3$. Khi đó $S_5$ bằng
	\choice 
		 {$11$}
		{$50$} 
		{$10$} 
		{\True  $25$}
	\loigiai{ 
		Ta có $S_n=nu_1+\dfrac{n\left(n-1\right)}{2}d \Rightarrow S_5=5u_1+\dfrac{5\cdot 4}{2}d=25$. 
	 }
\end{ex}

% nguồn: Ok_TracNghiemLop11_Lan1
\begin{ex}%[1D2H2-6]%[ma:27c43ba0]
	Cho cấp số cộng $(u_n)$ có  $u_1=3$ và $d=4$. Tổng $10$ số hạng đầu tiên của cấp số cộng đã cho là
	 \choice
		{$S_{10}=175$} 
		{$S_{10}=420$} 
		 {$S_{10}=230$}
		{\True  $S_{10}=210$}
	\loigiai{ 
		$S_n=\dfrac{n}{2}[2u_1+(n-1)d]\\ 
		\Rightarrow  S_{10}=\dfrac{10}{2}[2\cdot3+(10-1)\cdot4]=210$.
	 }
\end{ex}

% nguồn: Ok_Toan11_TracNghiem_Lan2
\begin{ex}%[1D2V2-6]%[ma:53d965b7]
	Cho một cấp số cộng $(u_n)$ có $u_1=1$ và tổng $100$ số hạng đầu bằng $10\,000$. Tính tổng  $S=\dfrac {1}{u_1u_2}+\dfrac {1}{u_2u_3}+\cdots+\dfrac {1}{u_{99}u_{100}}$.
	 \choice
		{$S=\dfrac  {100}{201}$}
		 {$S=\dfrac {200}{201}$}
		{$S=\dfrac {198}{199}$} 
		 {\True $S=\dfrac {99}{199}$}
	\loigiai{Gọi $d$ là công sai của cấp số cộng  đã cho.\\
		 Ta có
		$S_{100}=50(2u_1+99d)=10\,000\Rightarrow d=\dfrac  {200-2u_1}{99}=2$.
		 \allowdisplaybreaks
		\begin{align*}
			2S&=\dfrac {2}{u_1u_2}+\dfrac  {2}{u_2u_3}+\cdots +\dfrac {2}{u_{99}u_{100}}\\
			&=\dfrac  {u_2-u_1}{u_1u_2}+\dfrac {u_3-u_2}{u_2u_3}+\cdot \cdot \cdot +\dfrac {u_{99}-u_{100}}{u_{99}u_{100}}\\
			&=\dfrac  {1}{u_1}-\dfrac {1}{u_2}+\dfrac {1}{u_2}-\dfrac {1}{u_3}+\cdots +\dfrac {1}{u_{98}}-\dfrac {1}{u_{99}}+\dfrac {1}{u_{99}}-\dfrac {1}{u_{100}}\\
			&=\dfrac  {1}{u_1}-\dfrac {1}{u_{100}}\\
			&=\dfrac {1}{u_1}-\dfrac {1}{u_1+99d}\\ 
			&=\dfrac {198}{199}.\\ 
		\end{align*}
		Suy ra  $S=\dfrac {99}{199}$.
	} 
\end{ex}

% nguồn: Ok_Toan11_TracNghiem_Lan2
\begin{ex}%[1D2V2-6]%[ma:2557b065]
	Cho cấp số cộng $(u_n)$ có số hạng đầu bằng $1$ và tổng  $100$ số hạng đầu bằng $14\,950$. Giá trị của tổng $\dfrac {1}{u_1u_2}+\dfrac {1}{u_2u_3}+\cdots +\dfrac {1}{u_{49}u_{50}}$ bằng
	\choice 
		 {$\dfrac {49}{74}$}
		{$148$} 
		{\True $\dfrac {49}{148}$} 
		{$74$} 
	\loigiai{Gọi $ d$  là công sai của cấp số cộng. Ta có $S_{100}=50(2u_1+99d)=14\,950$ với $ u_1=1\Rightarrow d=3$.\\
		Đặt $S=\dfrac {1}{u_1u_2}+\dfrac {1}{u_2u_3}+\cdots +\dfrac {1}{u_{49}u_{50}}$.\\ 
		 Ta có
		 \allowdisplaybreaks
		\begin{align*}
			S\cdot d&=\dfrac  {d}{u_1u_2}+\dfrac {d}{u_2u_3}+\cdots +\dfrac {d}{u_{49}u_{50}}\\
			&=\dfrac {u_2-u_1}{u_1u_2}+\dfrac {u_3-u_2}{u_2u_3}+\cdots+\dfrac {u_{50}-u_{49}}{u_{49}u_{50}}\\ 
			&=\dfrac {1}{u_1}-\dfrac  {1}{u_{50}}\\
			 &=1-\dfrac {1}{1+49\cdot 3}\\
			&=\dfrac  {147}{148}.
		\end{align*}
		Với $d=3$ nên  $S=\dfrac {49}{148}$.
	} 
\end{ex}

% nguồn: Ok_Toan11_TracNghiem_Lan2
\begin{ex}%[1D2V2-6]%[ma:65f420d5]
	Cho  cấp số cộng $(u_n)$ có số hạng đầu bằng $1$ và tổng $100$ số hạng đầu bằng $14\,950$. Giá trị của tổng $\dfrac{1}{u_1u_2}+\dfrac{1}{u_2u_3}+\ldots+\dfrac{1}{u_{49}u_{50}}$ bằng
	 \choice
		{$\dfrac{49}{74}$} 
		 {$148$}
		{\True  $\dfrac{49}{148}$}
		{$74$} 
	 \loigiai{
		Gọi $d$ là công sai của cấp số cộng. Ta có $S_{100}=50(2u_1+99d)=14\,950$ với $u_1=1  \Rightarrow d=3$.\\
		 Đặt $S=\dfrac{1}{u_1u_2}+\dfrac{1}{u_2u_3}+\ldots+\dfrac{1}{u_{49}u_{50}}$.\\
		 Ta có
		\begin{eqnarray*}
			S\cdot d &=& \dfrac{d}{u_1u_2}+\dfrac{d}{u_2u_3}+\ldots+\dfrac{d}{u_{49}u_{50}} =  \dfrac{u_2-u_1}{u_1u_2}+\dfrac{u_3-u_2}{u_2u_3}+\ldots+\dfrac{u_{50}-u_{49}}{u_{49}u_{50}}\\
			&=&  \dfrac{1}{u_1}-\dfrac{1}{u_{50}}\\
			&=&1-\dfrac{1}{1+49\cdot  3}=\dfrac{147}{148}.
		\end{eqnarray*}
		Với $d=3$  nên $S=\dfrac{49}{148}$.
	} 
\end{ex}

% nguồn: Ok_Toan11_TracNghiem_Lan2
\begin{ex}%[1D2V2-6]%[ma:4fd10a84]
	Cho cấp số cộng $(u_n)$ biết $u_5=18$ và $4S_n=S_{2n}$. Giá trị $u_1$ và $d$ là 
	\choice 
		{$u_1=2;  d=3$}
		{$u_1=3;  d=2$}
		 {$u_1=2; d=2$}
		 {\True $u_1=2; d=4$}
	 \loigiai{
		Ta  có $u_5=18 \Leftrightarrow u_1+4d=18$.\\
		Lại  có
		\begin{eqnarray*}
			 &&4S_5=S_{10}\\
			& \Leftrightarrow& 4\left(5u_1+\dfrac{5\cdot 4}{2}d\right)=10u_1+\dfrac{10\cdot 9}{2}d\\  &\Leftrightarrow& 2u_1-d=0.
		\end{eqnarray*}
		Khi đó ta có  hệ phương trình $\heva{ &u_1+4d=18\\ &2u_1-d=0} \Leftrightarrow \heva{&u_1=2\\ &d=4.}$
	 }
\end{ex}

% nguồn: Ok_Toan11_TracNghiem_Lan2
\begin{ex}%[1D2V2-6]%[ma:0fd011f9]
	Cho tam giác đều $A_1B_1C_1$ có độ dài cạnh bằng $4$. Trung điểm của các cạnh tam giác $A_1B_1C_1$ tạo thành tam giác $A_2B_2C_2$, trung điểm của các cạnh tam giác $A_2B_2C_2$ tạo  thành tam giác $A_3B_3C_3$, $\cdots$ Gọi $P_1$, $P_2$, $P_3$, $\ldots$lần lượt là chu vi của tam giác $A_1B_1C_1$, $A_2B_2C_2$, $A_3B_3C_3$, $\ldots$. Tính tổng chu vi $P=P_1+P_2+P_3+\cdots$
	\choice 
		 {$P=8$}
		 {\True $P=24$}
		 {$P=6$}
		{$P=18$} 
	 \loigiai{Ta có
		$P_2=\dfrac {1}{2}P_1$; $P_3=\dfrac {1}{2}P_2=\dfrac {1}{4}P_1$; $P_4=\dfrac  {1}{2}P_3=\dfrac {1}{8}P_1 \cdots$; $P_n=\dfrac {1}{2^{n-1}}P_1\cdots$
		\allowdisplaybreaks 
		\begin{align*}
			P&=P_1+P_2+P_3+\cdots \\ 
			&=P_1+\dfrac  {1}{2}P_1+\dfrac {1}{4}P_1+\dfrac {1}{8}P_1+\cdots\\
			&=\dfrac {P_1}{1-\dfrac {1}{2}}\\ 
			&=2P_1\\ 
			&=24. 
		\end{align*}
		Vậy  $P=24$.
	 }
\end{ex}


%%%%%%%%%%%%%%%% ĐÚNG/SAI %%%%%%%%%%%%%%%%

% nguồn: Ok_dungsai_lop11_Lan2
\begin{ex}%[1D2N2-6]%[ma:f8669847]
	Cho cấp số cộng $(u_n)$ có số  hạng tổng quát là $u_n=3n-2$.
	\choiceTF 
		{\True Số hạng đầu của  cấp số cộng là $u_1=1$}
		{Công sai $d$ của cấp số cộng là  $d=-2$}
		{\True Tổng  của $20$ số hạng đầu của cấp số cộng là $S_{20}=590$}
		{Số $313$ là số hạng thứ $100$ trong  cấp số cộng}
	 \loigiai{
		\begin{itemchoice}
			\itemch {\bf Đúng}.\\ Vì $u_1=3\cdot 1-2=1$. 
			\itemch {\bf Sai}.\\ Vì  $d=u_{n+1}-u_n=3(n+1)-2-(3n-2)=3$.
			\itemch {\bf Đúng}.\\ Vì  $S_{20}=nu_1+\dfrac{n(n-1)}{2}\cdot d=20\cdot 1+\dfrac{20(20-1)}{2}\cdot 3=20+10\cdot 19\cdot 3=590$.
			 \itemch {\bf Sai}.\\ Vì
			\begin{eqnarray*}
				&  & u_n=u_1+(n-1)d\\
				&\Leftrightarrow & 313=1+(n-1)3\\ 
				&\Leftrightarrow &  313=3n-2\\
				&\Leftrightarrow  & 3n=315\\
				 &\Leftrightarrow & n=105.
			\end{eqnarray*}
		\end{itemchoice}
	} 
\end{ex}

% nguồn: Ok_dungsai_lop11_Lan2
\begin{ex}%[1D2N2-6]%[1D2H2-2]%[ma:7077c8c7]
	Cho cấp số cộng $(u_n)$ biết $u_1=1$; công sai $d$ và tổng của $100$  số hạng đầu tiên bằng $14\,950$.
	\choiceTF 
		 {\True Tổng $n$ số hạng đầu của cấp số cộng trên là $S_n=n+\dfrac{n(n-1)}{2}\cdot d$}
		{Công sai của  cấp số cộng trên là $d=2$}
		{Tổng của $15$ số hạng  đầu của cấp số cộng trên bằng $225$}
		{\True $S=\dfrac{1}{u_2\sqrt{u_1}+u_1\sqrt{u_2}}+\dfrac{1}{u_3\sqrt{u_2}+u_2\sqrt{u_3}}+\ldots  +\dfrac{1}{u_{2\,018}\sqrt{u_{2\,017}}+u_{2\,017}\sqrt{u_{2\,018}}}=\dfrac{1}{3}\left(1-\dfrac{1}{\sqrt{6\,052}}\right)$ }
	 \loigiai{
		\begin{itemchoice}
			\itemch {\bf Đúng}.\\ Ta có  $S_n=\dfrac{\left[2u_1+(n-1)d\right]n}{2}=\dfrac{2n+n(n-1)d}{2}=n+\dfrac{n(n-1)}{2}d$ (vì $u_1=1$).
			\itemch  {\bf Sai}.\\
			 Ta có
			\begin{eqnarray*}
				&  & S_{100}=100u_1+\dfrac{100\cdot 99}{2}d\\
				&\Leftrightarrow  & 100+4950d=14\,950\\
				&\Leftrightarrow & d=3. 
			\end{eqnarray*}
			\itemch {\bf Sai}.\\ Vì $S_{15}=15\cdot 1+\dfrac{15\cdot 14}{2}\cdot  3=330$.
			\itemch  {\bf Đúng}.\\ Vì $u_{2\,018}=u_1+2\,017d=6\,052$. Ta có
			\begin{eqnarray*}
				\dfrac{1}{u_{k+1}\sqrt{u_k}+u_k\sqrt{u_{k+1}}}&=&\dfrac{1}{\sqrt{u_k}\cdot  \sqrt{u_{k+1}}\cdot \left(\sqrt{u_k}+\sqrt{u_{k+1}}\right)}\\
				&=&\dfrac{1}{d}\cdot \dfrac{\sqrt{u_{k+1}}-\sqrt{u_k}}{\sqrt{u_k}\cdot \sqrt{u_{k+1}}}=\dfrac{1}{d}\cdot \left(\dfrac{1}{\sqrt{u_k}}-\dfrac{1}{\sqrt{u_{k+1}}}\right). 
			\end{eqnarray*}
			Suy  ra
			\begin{eqnarray*}
				S&=&\dfrac{1}{d}\cdot \left(\dfrac{1}{\sqrt{u_1}}-\dfrac{1}{\sqrt{u_2}}\right)+\dfrac{1}{d}\cdot  \left(\dfrac{1}{\sqrt{u_2}}-\dfrac{1}{\sqrt{u_3}}\right)+\ldots +\dfrac{1}{d}\cdot \left(\dfrac{1}{\sqrt{u_{2\,017}}}-\dfrac{1}{\sqrt{u_{2\,018}}}\right)\\
				&=&\dfrac{1}{d}\cdot \left(\dfrac{1}{\sqrt{u_1}}-\dfrac{1}{\sqrt{u_{2\,018}}}\right)\\ 
				&=&\dfrac{1}{3}\left(1-\dfrac{1}{\sqrt{6\,052}}\right). 
			\end{eqnarray*}
		\end{itemchoice}
	} 
\end{ex}

% nguồn: Ok_dungsai_lop11_Lan2
\begin{ex}%[1D2N2-6]%[ma:aae62dcb]
	Cho cấp số cộng $(u_n)$ có $7$ số hạng với $u_2=-1$  và công sai $d=-3$.
	 \choiceTF
		{Số hạng đầu của cấp số cộng này  là $u_1=-4$}
		{\True Số hạng thứ $7$ của cấp số  cộng này là $-16$}
		 {\True Tổng $7$ số hạng của cấp số cộng là $-49$}
		{Dạng khai triển của cấp số cộng là $2$; $-1$; $-4$; $-7$;  $-11$; $-13$; $-16$}
	\loigiai{ 
		\begin{itemchoice}
			\itemch {\bf Sai}.\\ Vì  $u_2=u_1+d\Rightarrow u_1=u_2-d=(-1)-(-3)=2$.
			\itemch {\bf Đúng}.\\ Ta dễ dàng liệt  kê các số hạng của dãy là $2$; $-1$; $-4$; $-7$; $-10$; $-13$; $-16$. Vậy $u_7=-16$.\\
			Hoặc áp dụng công thức số hạng tổng  quát $u_7=u_1+6d=-16$.
			\itemch {\bf Đúng}.\\ Liệt kê như câu b) rồi cộng các số hạng lại ta được $S=-49$.\\  Hoặc áp dụng công thức tổng $n$ số hạng đầu của một cấp số cộng $S=\dfrac{\left(u_1+u_n\right)n}{2}=-49$.
			 \itemch {\bf Sai}.\\ Theo câu b) ta được dãy là $2$; $-1$; $-4$; $-7$; $-10$; $-13$; $-16$. Vậy $u_5=-10\ne-11$.
		\end{itemchoice}
	} 
\end{ex}

% nguồn: Ok_dungsai_lop11_Lan2
\begin{ex}%[1D2N2-6]%[ma:447541b2]
	Cho cấp số  cộng $(u_n)$ có $u_1=123$, $u_3-u_{15}=84$.
	 \choiceTF
		 {\True Số hạng thứ $17$ của cấp số cộng là $u_{17}=11$}
		 {\True Công sai của cấp số cộng là $d=-7$}
		 {Số hạng thứ $2$ của cấp số cộng là $u_2=130$}
		{Tổng $17$ số hạng đầu tiên của cấp  số cộng là $S_{17}=1\,130$}
	 \loigiai{
		Giả sử cấp số cộng $(u_n)$ có công sai  $d$.\\
		Theo giả thiết ta có: 
		\begin{eqnarray*}
			& &  u_3-u_{15}=84\\
			&\Leftrightarrow & u_1+2d-u_1-14d=84\\ 
			&\Leftrightarrow & -12d=84\\ 
			&\Leftrightarrow  & d=-7.
		\end{eqnarray*}
		Vậy 
		\begin{itemize}
			 \item $u_{17}=u_1+16d=123+16(-7)=11$.
			 \item $u_2=u_1+d=123+(-7)=116$.
			 \item $S_{17}=\dfrac{\left(u_1+u_{17}\right)17}{2}=\dfrac{(123+11)17}{2}=1\,139$.
		\end{itemize}
		\begin{itemchoice}
			 \itemch {\bf Đúng}.\\
			\itemch {\bf Đúng}.\\ 
			\itemch {\bf Sai}.\\ 
			 \itemch {\bf Sai}.\\
		\end{itemchoice}
	} 
\end{ex}

% nguồn: Ok_dungsai_lop11_Lan2
\begin{ex}%[1D2N2-6]%[ma:267fcf90]
	Cho cấp số  cộng $(u_n)$ với $u_1=2$ và $u_2=6$.
	\choiceTF 
		{\True Công sai của cấp số cộng đã cho  là $d=4$}
		{Số hạng  thứ năm của cấp số cộng là $u_5=22$}
		{\True Tổng của $10$ số hạng đầu của cấp  số cộng là $S_{10}=200$}
		 {Số $130$ là số hạng thứ $30$ trong cấp số cộng}
	\loigiai{ 
		\begin{itemchoice}
			\itemch {\bf Đúng}.\\ Vì  $d=u_2-u_1=6-2=4$.
			\itemch {\bf Sai}.\\ Vì $u_5=u_1+4d=2+4\cdot 4=18$. 
			\itemch {\bf Đúng}.\\ Vì $S_{10}=nu_1+\dfrac{n(n-1)}{2}\cdot d=10\cdot 2+\dfrac{10(10-1)}{2}\cdot 4=20+5\cdot 9\cdot  4=200$.
			\itemch {\bf Sai}.\\  Vì
			\begin{eqnarray*}
				& & u_n=u_1+(n-1)d\\ 
				&\Leftrightarrow  & 130=2+(n-1)4\\
				&\Leftrightarrow & 130=4n-2\\ 
				&\Leftrightarrow &  4n=132\\
				&\Leftrightarrow  & n=33.
			\end{eqnarray*}
		\end{itemchoice}
	 }
\end{ex}

% nguồn: Ok_dungsai_lop11_Lan2
\begin{ex}%[1D2N2-6]%[ma:8b57962d]
	Cho dãy  cấp số cộng $(u_n)$ có $u_1=4$. Biết tổng $20$ số hạng đầu tiên bằng $460$.
	\choiceTF 
		{\True Dãy số $(u_n)$ có $d=2$} 
		{Dãy số $(u_n)$ có  $u_4=8$}
		{\True Dãy số $(u_n)$ có  $S_{10}=120$}
		 {\True Dãy số $(u_n)$ có hiệu $S_8-S_4=60$}
	\loigiai{ 
		Ta  có $S_{20}=\dfrac{20\left(u_1+u_{20}\right)}{2}=10(4+4+19d)=460\Leftrightarrow d=2$.\\
		Suy ra $\heva{&u_4=u_1+3d=4+3\cdot 2=10\\&S_{10}=\dfrac{10\left(u_1+u_9\right)}{2}=5(4+4+8\cdot  2)=120\\&S_8-S_4=\dfrac{8\left(u_1+u_8\right)}{2}-\dfrac{4\left(u_1+u_4\right)}{2}=4\left(2u_1+7d\right)-2\left(2u_1+3d\right)=4u_1+22d=60.}$\\
		\begin{itemchoice}
			\itemch {\bf Đúng}. 
			\itemch  {\bf Sai}.\\
			\itemch {\bf  Đúng}.\\
			 \itemch {\bf Đúng}.\\
		\end{itemchoice}
	} 
\end{ex}

% nguồn: Ok_dungsai_lop11_Lan2
\begin{ex}%[1D2N2-6]%[1D2N2-4]%[ma:b476d3c4]
	Cho cấp số cộng $(u_n)$ có  $u_4=-12$, $u_{14}=18$.
	 \choiceTF
		 {\True Công sai của cấp số cộng là $d=3$}
		{Số hạng đầu của cấp số cộng là  $u_1=21$}
		{\True  Số hạng thứ $9$ của cấp số cộng là $u_9=3$}
		{Tổng $5$ số hạng đầu của cấp số cộng  là $S_5=-60$}
	 \loigiai{
		Gọi  $u_1$ là số hạng đầu, $d$ là công sai của cấp số cộng $(u_n)$.\\
		Ta có 
		\begin{eqnarray*}
			& &  \heva{&u_4=u_1+3d \\&u_{14}=u_1+13d}\\
			 &\Leftrightarrow & \heva{&u_1+3d=-12\\&u_1+13d=18}\\
			 &\Leftrightarrow & \heva{&d=3\\&u_1=-21.}
		\end{eqnarray*}
		 Số hạng thứ $9$ là $u_9=u_1+8d=-21+8\cdot 3=3$.\\
		Tổng $5$ số hạng đầu  là $S_5=5u_1+\dfrac{5(5-1)}{2}d=-75$.
		\begin{itemchoice}
			\itemch  {\bf Đúng}.
			\itemch {\bf Sai}.\\ 
			\itemch {\bf Đúng}.\\ 
			\itemch {\bf Sai}.\\ 
		\end{itemchoice}
	} 
\end{ex}

% nguồn: Ok_dungsai_lop11_Lan1
\begin{ex}%[1D2H2-6]%[ma:34821178]
 Cho cấp số cộng $\left(u_n \right)$, biết rằng  $u_1=5$ và tổng của $50$ số hạng đầu bằng $5150$, khi đó
 \choiceTF 
  {Công sai của  cấp số cộng bằng $6$}
   {\True Số hạng $u_{85}=341$}
   {Số hạng $u_{10}=42$}
  {\True Tổng của 85 số hạng đầu $S_{85}=14705$} 
  \loigiai{
  \begin{itemchoice}
   \itemch Ta  có $S_{50}=\dfrac{50}{2} \left(2u_1+49d\right)=\dfrac{50}{2} (2\cdot 5+49d)=5150\Rightarrow d=4$.
   \itemch $u_n=u_1+(n-1)d=5+(n-1)4=1+4n$. Suy ra số hạng $u_{85}=341$. 
   \itemch $u_n=u_1+(n-1)d=5+(n-1)4=1+4n$.  Suy ra số hạng $u_{10}=41$.
   \itemch $S_{85}=\dfrac{85}{2} \left(2u_1+84d\right)=\dfrac{85}{2} (2\cdot  5+84\cdot 4)=14705$.
  \end{itemchoice}
 } 
\end{ex}

% nguồn: Ok_dungsai_lop11_Lan1
\begin{ex}%[1D2H2-6]%[ma:c3209c37]
 Cho cấp số cộng $\left(u_n \right)$, gọi $S_n$ là tổng $n$ số hạng đầu tiên của nó. Biết $S_7=77$ và $S_{12}=192$. Khi đó 
 \choiceTF 
  {\True Số hạng  $u_1=5$}
  {\True Tổng  $u_1+u_3=14$}
  {Công sai của cấp số cộng bằng $3$} 
  {\True  Số hạng $u_{11}=25$}
  \loigiai{
  \begin{itemchoice}
   \itemch Gọi  $d$ là công sai của cấp số cộng, ta có
   $\heva{&S_7=77\\&S_{12}=192}\Leftrightarrow \heva{&\dfrac{7}{2}\left( 2u_1+6d\right) =77\\&\dfrac{12}{2}\left( 2u_1+11d\right)  =192}\Leftrightarrow\heva{&7u_1+21d=77\\&12u_1+66d=192}\Leftrightarrow \heva{&u_1=5\\&d=2.}$
   \itemch  $u_1+u_3=2u_1+2d=14$.
    \itemch $d=2$.
   \itemch $u_{11}=u_1+10d=25$. 
  \end{itemchoice}
 } 
\end{ex}

% nguồn: Ok_dungsai_lop11_Lan1
\begin{ex}%[1D2H2-6]%[ma:7e5a56c8]
 Cho  cấp số cộng $\left(u_n\right)$ có số hạng đầu $u_1=-5$ và công sai $ d=3$.
 \choiceTF 
  {\True Số $100$ là số hạng thứ $36$ của cấp  số cộng}
  { Số hạng thứ $5$ của  cấp số cộng là $ 9$}
  {\True  Tổng $10$ số hạng đầu của cấp số cộng là $85$}
  {Số hạng tổng quát cảu cấp số cộng $\left(u_n\right)$ là  $\left(u_n\right)=3n-7$}
 \loigiai{ 
  \begin{itemchoice}
    \itemch Ta có
    $u_n=u_1+(n-1)d\Leftrightarrow 100=-5+(n-1)3\Leftrightarrow n=36$.
   \itemch Ta có $u_5=u_1+4d=-5+4\cdot3=7$. 
   \itemch Ta  có
   $S_n=\dfrac{n}{2}\left[2u_1+(n-1)d\right]\Rightarrow{S_{10}}=\dfrac{10}{2}\left[2\cdot(-5)+(10-1)3\right]=85$. 
   \itemch  $u_n=u_1+\left(n-1\right)d=-5+\left(n-1\right)3=3n-8$.
 \end{itemchoice}}
\end{ex}

% nguồn: Ok_dungsai_lop11_Lan1
\begin{ex}%[1D2H2-6]%[ma:49f98b7b]
 Cho cấp số cộng $\left(u_n\right)$có $u_n=5n-8$. Xác định tính  đúng, sai của các phát biểu sau đây
 \choiceTF 
  {\True $u_1=-3$} 
  {Công sai  $ d=-8$}
  {\True Số $492$ là  số hạng thứ $100$ của $\left(u_n\right)$}
  {Tổng của $10$ số hạng đầu tiên của cấp  số cộng bằng $ 190$}
 \loigiai{ 
  \begin{itemchoice}
   \itemch $u_1=5\cdot1-8=-3$ 
    \itemch $ d=u_2-u_1=2-\left(-3\right)=5$.
    \itemch $u_{100}=5\cdot100-8=492$.
    \itemch $S_{10}=\dfrac{10}{2}\left(2u_1+9d\right)=5\left[2\left(-3\right)+9\cdot5\right]=195$.
  \end{itemchoice}
 } 
\end{ex}

% nguồn: Ok_dungsai_lop11_Lan2
\begin{ex}%[1D2H2-6]%[ma:41b183f5]
	Cho cấp  số cộng $(u_n)$ có các số hạng đầu lần lượt là $5$; $9$; $13$; $17$; $\ldots$
	 \choiceTF
		{Công sai của cấp số cộng là $2$} 
		{\True Số hạng thứ sáu của cấp số cộng là  $25$}
		{\True Số hạng tổng quát $u_n$ của cấp  số cộng là $u_n=4n+1$}
		{Biết tổng $n$ số hạng đầu bằng  $230$. Khi đó cấp số cộng có $15$ số hạng}
	 \loigiai{
		\begin{itemchoice}
			\itemch  {\bf Sai}.\\ Vì $d=u_2-u_1=9-5=4$.
			\itemch {\bf  Đúng}.\\ Vì $u_6=u_1+5d=5+5\cdot 4=25$.
			\itemch {\bf Đúng}.\\ Vì  $u_n=u_1+(n-1)d=5+(n-1)\cdot 4=4n+1$.
			 \itemch {\bf Sai}.\\ Vì $230=S_n=nu_1+\dfrac{n(n-1)}{2}\cdot d=5n+\dfrac{n(n-1)}{2}\cdot 4=5n+2n(n-1)=2n^2+3n$.\\
			 Ta có phương trình
			\begin{eqnarray*}
				&&2n^2+3n=230\\ 
				&\Leftrightarrow& 2n^2+3n-230=0\\ 
				&\Leftrightarrow&  \hoac{&n=-\dfrac{23}{2}\, (\text{loại}) \\&n=10\, (\text{nhận}).}
			\end{eqnarray*}
			 Vậy cấp số cộng có $10$ số hạng.
	\end{itemchoice}}
\end{ex}

% nguồn: Ok_dungsai_lop11_Lan2
\begin{ex}%[1D2H2-6]%[ma:194d10e0]
	Một cấp số  cộng có năm số hạng mà tổng số hạng đầu và số hạng thứ tư bằng $36$, tổng của số hạng thứ hai và số hạng cuối bằng $44$. Xét tính đúng sai của các khẳng định sau
	\choiceTF 
		{Dãy cấp  số cộng có số hạng đầu $u_1=11$}
		{\True Dãy cấp số  cộng có tổng $u_1+u_5=40$}
		{\True Dãy cấp số cộng có  $u_2=16$}
		{Tổng của $3$ số hạng đầu tiên trong  dãy cấp số cộng bằng $45$}
	\loigiai{ 
		Ta  có
		\begin{eqnarray*}
			&&\heva {&u_1+u_4=36\\ &u_2+u_5=44} \\ &\Leftrightarrow& \heva{&u_1+(u_1+3d)=36\\&(u_1+d)+(u_1+4d)=44}\\& \Leftrightarrow&  \heva{&2u_1+3d=36\\ &2u_1+5d=44}\\
			&\Leftrightarrow& \heva{&u_1=12\\&  d=4.}\\
		\end{eqnarray*}
		Suy ra  $\heva{&u_2=u_1+d=16\\ &u_1+u_5=u_1+u_1+4d=40\\ &S_3=\dfrac{3(u_1+u_3)}{2}=48.}$
		\begin{itemchoice}
			\itemch  Dãy cấp số cộng có số hạng đầu $u_1=12$.
			\itemch  Dãy cấp số cộng có tổng $u_1+u_5=40$.
			\itemch Dãy cấp số cộng có $u_2=16$. 
			\itemch Tổng của  $3$ số hạng đầu tiên trong dãy cấp số cộng bằng $48$.
		\end{itemchoice}
	 }
\end{ex}

% nguồn: Ok_dungsai_lop11_Lan2
\begin{ex}%[1D2H2-6]%[ma:f3f60cee]
	Cho dãy cấp số cộng $(u_n)$ có $u_1=4$.  Biết tổng $20$ số hạng đầu tiên bằng $460$. Xét tính đúng sai của các khẳng định sau
	\choiceTF 
		{\True Dãy số $(u_n)$  có $d=2$}
		 {Dãy số $(u_n)$ có $u_4=8$}
		{\True Dãy số $(u_n)$ có  $S_{10}=120$}
		{\True  Dãy số $(u_n)$ có hiệu $S_8-S_4=60$}
	 \loigiai{
		 Ta có $S_{20}=\dfrac{20(u_1+u_{20})}{2}=10(4+4+19d)=460 \Leftrightarrow d=2$.\\
		 Suy ra
		\begin{itemize}
			\item $u_4=u_1+3d=4+3\cdot  2=10$.
			 \item $S_{10}=\dfrac{10(u_1+u_9)}{2}=5(4+4+8\cdot 2)=120$.
			 \item
			\begin{eqnarray*}
				 S_8-S_4&=&\dfrac{8(u_1+u_8)}{2}-\dfrac{4(u_1+u_4)}{2}\\
				 &=&4(2u_1+7d)-2(2u_1+3d)\\
				&=&4u_1+22d\\ 
				 &=&60.
			\end{eqnarray*}
		\end{itemize}
		\begin{itemchoice}
			 \itemch Dãy số $(u_n)$ có $d=2$.
			\itemch Dãy số $(u_n)$ có $u_4=8$. 
			\itemch Dãy số  $(u_n)$ có $S_{10}=120$.
			 \itemch Dãy số $(u_n)$ có hiệu $S_8-S_4=60$.
		\end{itemchoice}
	 }
\end{ex}

% nguồn: Ok_dungsai_lop11_Lan2
\begin{ex}%[1D2H2-6]%[ma:2d041c0e]
	
	Cho cấp số cộng $(u_n)$, gọi ${S_n}$ là  tổng $n$ số hạng đầu tiên của nó. Biết ${S_7=77}$ và ${S_{12}=192}$. Khi đó
	 \choiceTF
		{\True Số hạng ${u_1}=5$} 
		{\True Tổng ${u_1}+{u_3}=14$} 
		{Công sai của cấp số  cộng bằng $3$}
		{\True Số hạng  ${u_{11}}=25$}
	 \loigiai{
	Gọi $d$ là công sai của cấp số cộng.\\ 
	Ta có $\heva{& {S_7}=77 \\ & {S_{12}}=192 }$ $\Leftrightarrow $ $\heva{& \frac{7}{2}(  2{u_1}+6d)=77 \\ & \frac{12}{2}( 2{u_1}+11d)=192 }$ $\Leftrightarrow$ $\heva{& 7{u_1}+21d=77 \\ & 12{u_1}+66d=192 }$ $\Leftrightarrow$ $\heva{& {u_1}=5 \\ & d=2 }$.\\
	Khi đó  ${u_n=u_1+(n-1) d=5+2(n-1)=3+2 n}$.
	\begin{itemchoice}
		\itemch \textbf{Đúng.} 
		\itemch  \textbf{Đúng.}
		\itemch  \textbf{Sai.}\\
		${u_1}=5$,  $u_{3}=3+2\cdot3=9$ nên ${u_1}+{u_3}=14$.
		\itemch  \textbf{Đúng.}\\
		$u_{11}=3+2\cdot11=25$. 
	\end{itemchoice}
	 }
\end{ex}


%%%%%%%%%%%%%%%% TRẢ LỜI NGẮN %%%%%%%%%%%%%%%%

% nguồn: Ok_TraLoiNgan_Lop11_Lan1
\begin{ex}%[1D2H2-6]%[ma:92177595]
	 Giải phương trình sau $(2x+1)+(2x+6)+(2x+11)+\cdots+(2x+96)=1010$.
	 \shortans{1}
	\loigiai{ 
		Ta có dãy số $(2x+1)$, $(2x+6)$, $(2x+11)$,$\ldots, (2x+96)$ lập thành cấp  số cộng có $\heva{&{u_1=2x+1} \\
		&{d=5} \\ 
		&{u_n=2x+96} \\ 
		 &{S_n=1010}.
		}$\\ 
		Suy ra  $2x+96=2x+1+(n-1)5\Rightarrow n=20$.\\
		Vậy $S_{20}=1010\Leftrightarrow \dfrac{20}{2}  \left(u_1+u_n \right)=1010\Leftrightarrow 10(2x+1+2x+96)=1010\Rightarrow x=1$.
	} 
\end{ex}

% nguồn: Ok_TraLoiNgan_Lop11_Lan2
\begin{ex}%[1D2H2-6]%[ma:7aa34217]
	Cho cấp số cộng  $(u_n)$ có ${u_{27}}+u_2=83$. Khi đó tổng $28$ số hạng đầu tiên của cấp số cộng $(u_n)$ là
	\shortans{1,16} 
	 \loigiai{
		Gọi $d$ và $u_1$ lần lượt là công sai và số  hạng đầu của cấp số cộng $(u_n)$.\\
		 Ta có
		\[{S_{28}}=\dfrac{28\left(u_1+{u_{28}}\right)}{2}=\dfrac{28\left(u_2-d+{u_{27}}+d\right)}{2}=\dfrac{28\left(u_2+{u_{27}}\right)}{2}=\dfrac{28\cdot  83}{2}=1\,162.\]
	 }
\end{ex}

% nguồn: Ok_TraLoiNgan_Lop11_Lan2
\begin{ex}%[1D2H2-6]%[ma:de101e57]
	 Chu vi của một đa giác là $158$ cm, số đo các cạnh của nó lập thành một cấp số cộng với công sai $d=3$ cm Biết cạnh lớn nhất là $44$ cm Khi đó số cạnh của đa giác đó là
	\shortans{4} 
	 \loigiai{
		Gọi cạnh nhỏ nhất là $u_1$ và số cạnh  của đa giác là $n$ với $n\in \mathbb{N^*}$.\\
		Ta có $44=u_1+(n-1)\cdot 3\Leftrightarrow  u_1=47-3n$.\\
		Vì chu vi của tứ giác là $158$ nên ta có 
		\[158=\dfrac{n(44+47-3n)}{2}\Leftrightarrow 3n^2-91n+316=0\Leftrightarrow \hoac{&n=\dfrac{79}{3}\\&n=4.}\] 
		 Vì $n\in \mathbb{N^*}$ nên $n=4$.
	 }
\end{ex}

% nguồn: Ok_TraLoiNgan_Lop11_Lan2
\begin{ex}%[1D2H2-6]%[ma:e0e21c17]
	Cho cấp số cộng $(u_n)$ với  $u_2=3$ và $u_5=12$. Giá trị $759$ là tổng của bao nhiêu số hạng đầu của cấp số cộng?
	\shortans{23} 
	 \loigiai{
		Ta  có
		$\heva{&u_2=u_1+d=3\\&u_5=u_1+4d=12}\Leftrightarrow  \heva{&u_1=0\\&d=3}$.\\
		Gọi tổng của $n$,  $(n>0, n\in\mathbb{N})$ số hạng đầu bằng $759$ suy ra
		\begin{eqnarray*}
			 &&S_n=\dfrac{n\left[2u_1+(n-1)d\right]}{2}=759\\
			&\Leftrightarrow& \dfrac{n[0+(n-1)\cdot  3]}{2}=759\\
			 &\Leftrightarrow& n(n-1)=506\\
			&\Leftrightarrow&  \hoac{&n=-22\quad \text{(loại)}\\&n=23.}
		\end{eqnarray*}
	} 
\end{ex}

% nguồn: Ok_TraLoiNgan_Lop11_Lan2
\begin{ex}%[1D2H2-6]%[ma:54433942]
	Một cấp số cộng có số hạng đầu bằng $5$ và công sai bằng $2$. Hỏi phải lấy tồng của bao nhiêu số hạng đầu của  cấp số cộng này để có tồng bằng $2\,700$?
	 \shortans{50}
	\loigiai{ 
		Giải sử phải lấy $n$, $(n\in \mathbb{Z^+})$ số hạng đầu tiên của cấp số cộng  để có tổng bằng $2\,700$.\\
		Ta có 
		\begin{eqnarray*}
			 &&S_n=\dfrac{\left(2u_1+(n-1)d\right)n}{2}\\
			 &\Leftrightarrow& \dfrac{(10+2(n-1))n}{2}=2\,700\\
			 &\Leftrightarrow& n^2+4n-2\,700=0\\
			 &\Leftrightarrow& n=50.
		\end{eqnarray*}
		Vậy phải lấy $50$  số hạng đầu của cấp số cộng.
	 }
\end{ex}

% nguồn: Ok_TraLoiNgan_Lop11_Lan2
\begin{ex}%[1D2H2-6]%[ma:609bbd2d]
	Cho cấp số cộng $(u_n)$ có $u_5=-15$, ${u_{20}}=60$. Tính tổng $10$ số hạng đầu tiên  của cấp số cộng đó.
	 \shortans{-125}
	 \loigiai{
		Gọi $u_1$, $d$ lần lượt là số hạng đầu và  công sai của cấp số cộng.\\
		Ta có $\heva{&u_5=-15\\&{u_{20}}=60}\Leftrightarrow \heva{&u_1+4d=-15\\&u_1+19d=60}\Leftrightarrow \heva{&u_1=-35\\&d=5.}$\\ 
		Vậy $S_{10}=\dfrac{10}{2}\cdot (2u_1+9d)=5\cdot \left[2\cdot  (-35)+9\cdot 5\right]=-125$.
	} 
\end{ex}

% nguồn: Ok_TraLoiNgan_Lop11_Lan2
\begin{ex}%[1D2H2-6]%[ma:a67335f3]
	Cho cấp số cộng $(u_n)$ có $u_3+u_{28}=100$. Hãy tính tổng của $30$ số hạng đầu  tiên của cấp số cộng đó.
	\shortans{1,50} 
	 \loigiai{
		Ta có  $u_3+u_{28}=100 \Leftrightarrow (u_1+2d)+(u_1+27d)=100 \Leftrightarrow 2u_1+29d=100$\\
		Mà $S_{30}=\dfrac{30}{2}(2u_1+29d)$ nên $S_{30}=15\cdot  100=1\,500$.
	 }
\end{ex}

% nguồn: Ok_TraLoiNgan_Lop11_Lan2
\begin{ex}%[1D2H2-6]%[ma:86dd0e92]
	Cho dãy số $(u_n)$ xác định bởi $\heva{&u_1=2\\& u_{n+1}=u_n+3, n \ge 1}$. Tìm  tổng $20$ số hạng đầu của dãy số $(u_n)$.
	\shortans{610} 
	\loigiai{ 
		Vì dãy số $(u_n)$ có $u_{n+1}=u_n+3, n \ge 1$ nên $(u_n)$ là cấp số cộng có $u_1=2$, công  sai $d=3$.\\
		Vậy tổng $20$ số hạng đầu của dãy số $(u_n)$  là $S_{20}=\dfrac{20}{2}[2u_1+19d]=10(2\cdot 2+19\cdot 3)=610$.
	} 
\end{ex}

% nguồn: Ok_TraLoiNgan_Lop11_Lan1
\begin{ex}%[1D2V2-6]%[ma:2185af08]
	Vào năm $2023$, dân số của một thành phố là khoảng $1{,}5$ triệu người. Giả sử mỗi năm, dân số của thành phố này tăng thêm khoảng $23$ nghìn người. Hãy ước tính dân số của thành phố này vào năm $2028$ bằng bao  nhiêu nghìn người? \shortans{1615}
	\loigiai{Dân số mỗi năm của thành phố từ năm $2023$ đến năm $2028$ lập thành cấp số cộng gồm $6$  số hạng với $u_1 = 1\,500\,000$ và công sai $d = 23\,000$.\\
		Theo đề bài ta có $$u_{6} = u_1 +  (6 - 1)d = 1\,500\,000 + 5 \cdot 23\,000 = 1\,615\,000.$$
	Vậy đến năm  $2028$ thì dân số khoảng $1\,615$ nghìn người.}
\end{ex}

% nguồn: Ok_TraLoiNgan_Lop11_Lan1
\begin{ex}%[1D2V2-6]%[ma:23ebdf12]
	Cho  hai cấp số cộng $(u_n)$ và $(v_n)$ có tổng của $n$ số hạng đầu tiên lần lượt là $S_n$, $T_n$. Biết $\dfrac{S_n}{T_n}=\dfrac{4n+1}{3n+1}$ với mọi $n\in\mathbb{N}^*$, tính $\dfrac{u_{17}}{v_{17}}$.\\
	 \shortans{1,33}
	\loigiai{ 
		Thay $n=33$ ta  có $\dfrac{S_{33}}{T_{33}}=\dfrac{133}{100}$.\\
		Mà $S_{33}=\dfrac{(u_1+u_{33})\cdot 33}{2}$, $T_{33}=\dfrac{(v_1+v_{33})\cdot 33}{2}$.\\ 
		Vậy $\dfrac{133}{100}=\dfrac{u_1+u_{33}}{v_1+v_{33}}=\dfrac{2u_{17}}{2v_{17}}=\dfrac{u_{17}}{v_{17}}$. 
		Từ đó suy ra  $\dfrac{u_{17}}{v_{17}}=1{,}33$.
	} 
\end{ex}

% nguồn: Ok_TraLoiNgan_Lop11_Lan1
\begin{ex}%[1D2V2-6]%[ma:9104d92d]
	 Tìm số hạng đầu tiên, công sai và tổng của $20$ số hạng đầu tiên của cấp số cộng $(u_n)$ biết rằng $\heva{& {u_5}=19 \\
	& {u_9}=35}$. 
	\shortans{820} 
	 \loigiai{
		Áp dụng công thức  $u_n=u_1+(n-1)d$, ta có $\heva{
		& {u_5}=19 \\ 
		& {u_9}=35  } \Leftrightarrow \heva{
		 & {u_1}+4d=19 \\
		& {u_1}+8d=35 }\Leftrightarrow \heva{ 
		& {u_1}=3  \\
		&  d=4.}$\\
		Vậy số hạng đầu tiên $u_1=3$, công sai  $d=4$.\\
	Tổng của $20$ số hạng đầu tiên: ${S_{20}}=\dfrac{20\left( 2u_1+19d  \right)}{2}=10\left( 2\cdot3+19\cdot4 \right)=820$.}
\end{ex}

% nguồn: Ok_TraLoiNgan_Lop11_Lan1
\begin{ex}%[1D2V2-6]%[ma:50c4d41b]
	Tìm tổng của $20$ số hạng đầu tiên của cấp số cộng $\left( u_n\right)$, biết rằng $$\heva{&u_2-u_3+u_5=10 \\&  u_2+u_7=17.}$$
	\shortans{-190} 
	 \loigiai{
		Áp dụng công thức $u_n=u_1+(n-1)\cdot  d$.\\
		Ta  có:
		 \allowdisplaybreaks
		\begin{eqnarray*}
			&&\heva{&u_2-u_3+u_5=10 \\& u_2+u_7=17}\\ &\Leftrightarrow &\heva{& (u_1+d)-(u_1+2d)+(u_1+4d)=10  \\& (u_1+d)+(u_1+6d)=17}\\
			&\Leftrightarrow  &\heva{& u_1+3d=10 \\& 2u_1+7d=17}\\ &\Leftrightarrow& \heva{&u_1=19\\& d=-3.}
		\end{eqnarray*}
		Khi đó $u_{20}=u_1+19d=19+19\cdot (-3)=-38$  và
		$S_n=\dfrac{n}{2}(u_1+u_n)$.\\ 
		$\Rightarrow  S_{20}=\dfrac{20}{2}\left[19+(-38) \right]=-190$.
	} 
\end{ex}

% nguồn: Ok_TraLoiNgan_Lop11_Lan1
\begin{ex}%[1D2V2-6]%[ma:6e4ce1b1]
	Một khán đài của một Nhà thi đấu thể thao được thiết kế với $20$ hàng ghế, trong đó hàng thứ nhất có $25$ ghế ngồi; hàng thứ hai có $28$ ghế ngồi,  hàng thứ ba là $31$ ghế ngồi,$\ldots$ và cứ như vậy cho đến hàng cuối cùng (số ghế ở hàng sau nhiều hơn $3$ ghế so với số ghế ở hàng liền trước nó). Tính tổng số ghế ngồi của khán đài đó.
	 \shortans{1070}
	\loigiai{ 
		\begin{itemize}
			\item Gọi  số ghế ở hàng thứ nhất là $u_1$ thì $u_1=25$.
			 \item Gọi số ghế ở hàng thứ hai là $u_2$ thì $u_2=28$.
			\item Gọi số ghế ở hàng thứ  ba là $u_3$ thì $u_3=31$.
			\item $\cdots$ 
			\item Gọi số ghế ở hàng thứ hai  mươi là $u_{20}$.
			\item Ta có dãy số $u_1$, $u_2$, $\cdots$, $u_{20}$ là cấp số cộng có  công sai $d=3$.
		\end{itemize}
		Tổng số ghế ngồi của khán đài là $S_{20}=\dfrac{20}{2}\left( 2\cdot u_1+19\cdot d\right) =10(50+57)  =1070$ ghế.
	 }
\end{ex}


%%%%%%%%%%%%%%%% TỰ LUẬN %%%%%%%%%%%%%%%%

% nguồn: Ok_TuanLuan_Toan11_Lan2
\begin{ex}%[1D2H2-6]%[ma:d5b9fc73]
	Tính tổng  $S=1-2+3-4+5-6+\ldots + (2n-1)-2n$ với $n \ge 1$ và $n \in \mathbb{N}$.
	\loigiai{ 
		Với mọi $n \in  \mathbb{N}^*$ thì $(2n-1)-2n = -1$.\\
		Ta  có $S=(1-2)+(3-4)+(5-6)+ \ldots+((2n-1)-2n)$.\\
		Do đó ta xem $S$ là tổng của $n$ số hạng, mà mỗi số hạng đều  bằng $-1$ nên $S = -n$.
	} 
\end{ex}

% nguồn: Ok_TuanLuan_Toan11_Lan2
\begin{ex}%[1D2H2-6]%[ma:e16bfc0a]
	Cho một cấp số cộng $(u_n)$ có $u_1 = 5$ và tổng của $50$ số hạng đầu bằng $5\,150$. Tìm công thức của số  hạng tổng quát $u_n$.
	\loigiai{ 
		Ta có $S_{50} = \dfrac{50}{2}\cdot (2u_1+49\cdot  d) = 5150 \Rightarrow d = 4$.\\
		Số hạng tổng quát của cấp số cộng bằng $u_n = u_1+(n-1)\cdot d  = 1+4n$.
	} 
\end{ex}

% nguồn: Ok_TuanLuan_Toan11_Lan2
\begin{ex}%[1D2H2-6]%[ma:f289751a]
	Chu vi một đa giác  là $158$, số đo các cạnh của nó lập thành một cấp số cộng với công sai $d=3$. Biết cạnh lớn nhất là $44$. Số cạnh của đa giác đó là
	\loigiai{ 
		Giả  sử đa giác đã cho có $n$ cạnh thì chu vi của đa giác là $S_n = \dfrac{(u_1+u_n)\cdot n}{2}$ với $u_1$ là cạnh nhỏ nhất.\\
		Suy ra $158 = \dfrac{(u_1+44)\cdot n}{2} \Leftrightarrow 316 = (u_1+44)\cdot n \Leftrightarrow 2^2\cdot 79 = (u_1+44)\cdot  n$.\\
		Do đó $u_1+44$ là ước nguyên dương của $316 = 2^2\cdot 79$ và đa giác có ít nhất ba  cạnh nên $\dfrac{316}{3} > u_1+44 > 44$.\\
		Suy ra $u_1+44 = 79 \Leftrightarrow u_1 =  35$.\\
		Số cạnh của đa  giác đã cho là $\dfrac{44-35}{3}+1 = 4$ (cạnh).
	 }
\end{ex}

% nguồn: Ok_TuanLuan_Toan11_Lan2
\begin{ex}%[1D2H2-6]%[ma:d0cea026]
	Cho cấp số cộng $\left(u_{n}\right)$ có $u_{5}=-15$, $u_{20}=60$. Tính tổng của $10$ số hạng đầu tiên của cấp số cộng đã cho. 
	\loigiai{Gọi $u_{1}$, $d$  lần lượt là số hạng đầu và công sai của cấp số cộng.\\
		Ta có  $\heva{&u_5=-15\\
		 &u_{20}=60}
		\Leftrightarrow 
		\heva{&u_1+4d=-15\\ 
		 &u_1 +19d=60
		 \Leftrightarrow}
		 \heva{&u_1=-35\\
	&d=5.}$}\\ 
	Vậy $S_{10}=\dfrac{10}{2} \cdot\left(2 u_{1}+9 d\right)=5 \cdot[2 \cdot(-35)+9  \cdot 5]=-125$.
\end{ex}

% nguồn: Ok_TuanLuan_Toan11_Lan2
\begin{ex}%[1D2H2-6]%[ma:57ebeaf7]
	Cho  cấp số cộng $\left(u_{n}\right)$ có $u_{4}=-12$, $u_{14}=18$. Tính tổng $16$ số hạng đầu tiên của cấp số cộng này.
	\loigiai{Gọi $d$ là công sai của cấp số cộng. Theo giả thiết,  ta có
		 $\heva{&u_1+3d=-12\\
		 &u_1+13d=18}\Leftrightarrow
		\heva{&u_1=-21\\ 
		 &d=3.}$\\
		Khi đó $S_{16}=\dfrac{\left(2 u_{1}+15 d\right) \cdot 16}{2}=8(-42+45)=24$. 
	} 
\end{ex}

% nguồn: Ok_TuanLuan_Toan11_Lan2
\begin{ex}%[1D2H2-6]%[ma:7b8fb3a7]
	Cho cấp số cộng có $u_{1}=1$ và công sai $d=-2$. Tổng $n$ số hạng  đầu tiên của cấp số cộng này là $S_{n}=-9\,800$. Tìm $n$.
	\loigiai{ 
	Ta  có $S_{n}=\dfrac{n}{2}\left[2 u_{1}+(n-1) d\right]=-9\,800 \Leftrightarrow n[2-2(n-1)]+19\,600=0 \Leftrightarrow n=100$.}
\end{ex}

% nguồn: Ok_TuanLuan_Toan11_Lan2
\begin{ex}%[1D2V2-6]%[ma:0fa1feda]
	Cho cấp số cộng $\left(u_{n}\right)$ biết $u_{5}=18$ và $4 S_{n}=S_{2 n}$. Tìm số hạng đầu tiên $u_{1}$ và công  sai $d$ của cấp số cộng.
	\loigiai{Ta có $u_{5}=18 \Leftrightarrow u_{1}+4 d=18 \quad (1)$.\\ 
		Ta  có
		\begin{eqnarray*}
			& &4 S_{n}=S_{2  n}\\
			&\Leftrightarrow &4\left[n u_{1}+\dfrac{n(n-1)  d}{2}\right]=\left[2 n u_{1}+\dfrac{2 n(2 n-1) d}{2}\right]\\
			&\Leftrightarrow &u_{1}+2  n d-2 d=2 u_{1}+2 n d-d\\
			&\Leftrightarrow &2 u_{1}-d=0 \quad(2). 
	\end{eqnarray*}}
	Từ (1) và (2) suy ra $u_{1}=2$ và $d=4$. 
\end{ex}
